\chapter{Simultaneous Localization and Mapping~(SLAM) in Dynamic Environments}\label{chap:slam}

This chapter presents RicoSLAM, a graph-based 2D~\gls{slam} framework for dynamic
environments developed within the scope of this PhD~Thesis, building on the 2D
point cloud registration framework introduced in \chapref{chap:distance-maps}.
RicoSLAM integrates the distance map-based tracker from
\chapref{chap:distance-maps} as its front-end with the
srrg2\_solver~\cite{grisetti:robotics:2020} optimization back-end to construct
an $SE(2)$ pose graph, detecting loop closures by combining nodes' proximity
in the world frame with graph-distance of the loop chain.
The proposed \gls{slam} framework filters dynamic observations
online during local map maintenance using a 2D laser scanner ray model,
while providing an offline 2D occupancy grid map generation step in which
the same ray model is combined with
Bresenham-like~\cite{bresenham:ibm:1965,zingl:uas:2016} line rasterization to
ignore dynamic observations when projecting all scans considered during online
\gls{slam} onto the final grid map. Additionally, this chapter presents
novel formulations for \gls{gba} based on distance map-based per-point factors,
building on the point-to-point and point-to-plane distance map-based
formulations presented in \chapref{chap:distance-maps}.
Experimental results indicate that RicoSLAM can operate in highly dynamic,
noisy environments while filtering dynamic observations from
the 2D~occupancy grid representation of the environment.

This chapter is organized as follows.
First, \secref{sec:slam:ricoslam} presents a high-level overview of the RicoSLAM
framework, positioning the contributions of this chapter in relation to
the components inherited from \chapref{chap:distance-maps}.
Next, \secref{sec:slam:online} describes the online \gls{slam} pipeline,
covering the $SE(2)$ pose-pose geodesic error factor and its Jacobians via
right-sided perturbation available in the
srrg2\_solver~\cite{grisetti:robotics:2020} framework, the local map splitting
and merging criteria, the loop closure detection procedure, the relocalization
mechanism, and the pose history streams (laser odometry, online and offline
corrected poses) logged for trajectory evaluation.
\secref{sec:slam:dynamics} describes the dynamics removal strategy proposed in
this work, both online (during local map maintenance) and offline
(during 2D~occupancy grid map generation).
Moreover, \secref{sec:slam:gba} introduces the distance map-based \gls{gba}
per-point formulations proposed in this work, deriving the
point-to-point and point-to-plane error factors together with their Jacobians
via right-sided perturbation on the $SE(2)$ manifold.
Then, \secref{sec:slam:results} presents the experimental evaluation, including
a benchmark on the IILABS~3D~\cite{ribeiro:access:2025} dataset in terms of
accuracy and scalability, the impact of dynamics removal from the 2D occupancy
grid representation, and a synthetic benchmark to evaluate the distance
map-based \gls{gba} proposed in this work.
Lastly, \secref{sec:slam:conclusions} summarizes the main conclusions and
directions for future work.

\section{Framework Overview}\label{sec:slam:ricoslam}

\begin{figure}[t]
  \centering
  \includegraphics[width=\textwidth]{slam/system/system-overview.drawio.pdf}
  \caption[Block diagram of the RicoSLAM framework]{%
  Block diagram of the RicoSLAM framework.}
  \label{fig:slam:ricoslam}
\end{figure}

RicoSLAM is a 2D~graph \gls{slam} framework compatible with 2D~laser scanners
and wheeled odometry for dynamic environments in industrial \gls{amr} applications.
This framework represents the robot trajectory and the environment map as
a sparse factor graph~\cite{grisetti:itsm:2010,dellaert-kaess:ftr:2017}.
Each $SE(2)$ variable in the graph encodes the pose of a keyframe in
a common world frame, together with its measurement:
the locally built point cloud, the precomputed distance map representation
introduced in \secref{sec:distance-maps:alignment:dmaps}, and
the 2D~laser scanner messages aggregated within that keyframe. Each factor
encodes a relative-pose measurement between two keyframes, generated either by
odometric chaining of the front-end tracker or by loop closure detection.

\figref{fig:slam:ricoslam} presents a block diagram of RicoSLAM.
Internally, the system is organized around two top-level containers -- a tracker
front-end and a factor graph back-end -- with two additional offline
post-processing stages, a distance map-based \gls{gba} procedure and a
2D~occupancy grid mapper, that operate on the optimized pose graph once the
online run has completed. All the blocks in the diagram are summarized as follows:
\begin{itemize}
\item The tracker is built on top of the distance map-based aligner from
      \chapref{chap:distance-maps} and is extended in this chapter with a
      dynamics-aware local mapping stage. Preprocessing
      (conversion of laser scan messages into a point cloud representation,
      voxelization, and normal computation) and the distance map-based
      point-to-point and point-to-plane formulations
      (see~\secref{sec:distance-maps:alignment:pt2pt} and
      \secref{sec:distance-maps:alignment:pt2pl}) are inherited from
      \chapref{chap:distance-maps}. We also added to the framework both
      point-to-point~\cite{besl-mckay:tpami:1992} and
      point-to-plane~\cite{chen-medioni:icra:1991} \gls{icp} formulations for
      2D~point cloud registration with the tracker's aligner,
      using distance map-guided data association.
      New to this chapter are the split and merge triggering criteria,
      which monitor the aligner output stream
      (translation, orientation, viewpoint, inlier ratio, and the
      timestamp from the latest laser scan message processed by the tracker).
      The split criterion determines when to commit the current keyframe to
      the graph and start tracking on a new keyframe.
      The merge criterion determines when to merge new returns into
      the tracker's local map, using a merger container that performs
      online ray-based dynamics removal on that local map by treating
      each laser beam as a sensor model that may see through a previously
      mapped point or add new observations to the local representation.
      The tracker maintains the current keyframe --
      local point cloud, precomputed distance map, and array of raw laser scans
      aggregated since the last split -- and exposes its
      laser odometry estimate and alignment statistics to the back-end;

\item The factor graph back-end uses srrg2\_solver~\cite{grisetti:robotics:2020}
      to solve the optimization problem and host all the components that operate
      on the graph. Its core data structure is the pose graph,
      populated by the tracker through the update-on-split action that
      promotes the current keyframe to a graph variable and adds a geodesic
      factor against the previously committed keyframe.
      The loop finder and graph visitor containers traverse the graph
      in search of loop closure candidates by combining world-frame proximity
      around the current robot pose with a graph-distance constraint that
      ensures non-trivial loop geometry. Each candidate is validated by
      the loop aligner, aligning the latest point cloud processed by the tracker
      against the candidate keyframe's local map, using two complementary
      initial guesses (pose graph-based and same-place-based). This loop aligner
      may be the same as the tracker's front-end or a different one.
      After a successful loop closure, a relative-pose geosedic factor is added
      between the current keyframe and the loop variables. Also, if
      the relocalization policy is enabled, it updates the tracker's keyframe
      from an existing one on the graph, relocalizing the robot in
      the existing back-end graph. The Gauss-Newton solver, provided by
      srrg2\_solver~\cite{grisetti:robotics:2020} and the only third-party
      component of the framework, performs iterative least-squares optimization
      on the $SE(2)$ manifold with a sparse-block Cholesky linear solver,
      and is invoked incrementally whenever a loop closure factor is added to
      the graph. Furthermore, a pose history container makes available the three
      pose streams produced over time: the raw laser odometry from the tracker,
      the online \gls{slam} pose estimate, and the corrected pose based on
      the online estimate while considering the loop closures throughout time.
      RicoSLAM generates TUM~\cite{sturm:iros:2012} files for these three pose
      streams through external \gls{ros} services upon user request.
      Additionally, RicoSLAM saves the factor graph through the serialization
      framework made available through the
      srrg2\_solver~\cite{grisetti:robotics:2020} into a human-readable JSON
      file, including keyframe variables and their global poses,
      the $SE(2)$ geosedic pose-pose factors, 2D~point clouds, and
      the arrays of 2D~laser scan messages associated with each keyframe;

\item Finally, two offline post-processing stages consume the optimized pose
      graph from the serialized JSON file after the online run completes.
      The \gls{gba} replaces the pose-pose constraints obtained from
      the tracker and the loop aligner by per-scan-point distance map residuals
      between different laser scan messages, thereby refining all scans that
      contributed to the online \gls{slam} map through novel distance map-based
      \gls{gba} formulations proposed in this chapter.
      The 2D~occupancy grid mapper generates a ROS-compatible occupancy grid map
      by Bresenham-based~\cite{bresenham:ibm:1965,zingl:uas:2016} rasterization
      of each retained scan, using the same laser ray model as
      the online merger to ignore the contributions of dynamic observations.
\end{itemize}

Throughout this chapter, the notation introduced in
\secref{sec:distance-maps:alignment:problem} is kept whenever possible.
Consistent with the convention adopted in \chapref{chap:distance-maps}
for $SE(2)$ optimization variables, a keyframe variable in the pose graph
is denoted by $\mathbf{x} \in SE(2)$ when it refers to a generic graph node,
and is given subscripts \emph{from} or \emph{to} when both endpoints of
a pose-pose geosedic factor must be distinguished.
The current pose \gls{wrt} the world frame is denoted
$\mathbf{x}_{R,W} = \mathbf{x}_{\mathrm{kf}} \cdot \mathbf{x}_{R,\mathrm{kf}}$,
where $\mathbf{x}_{\mathrm{kf}}$ is the current keyframe estimate and
$\mathbf{x}_{R,\mathrm{kf}}$ is the local laser odometry estimate of
the robot \gls{wrt} that keyframe. The right-sided perturbation update
$\mathbf{x} \oplus \boldsymbol{\xi} = \mathbf{x}\cdot\exp(\boldsymbol{\xi})$
adopted in
\secref{sec:distance-maps:background:least-squares:manifold} and
\secref{sec:distance-maps:alignment:problem} is used in all error and
Jacobian derivations of this chapter.

\section{Online SLAM}\label{sec:slam:online}

TBC

\section{Dynamics Removal}\label{sec:slam:dynamics}

TBC

\section{Global Bundle Adjustment}\label{sec:slam:gba}

TBC

\section{Experimental Results and Discussion}\label{sec:slam:results}

\subsection{Trajectory accuracy on the IILABS3D Dataset}\label{sec:slam:results:iilabs3d}

\todo{TBC}

{\setlength{\tabcolsep}{2pt}
\begin{table}[!t]
\centering
\caption[%
    Trajectory accuracy of RicoSLAM, SLAM~Toolbox, and GMapping on
    four planar sequences of the IILABS~3D dataset
    (Nav~A~Diff, Nav~A~Omni, Loop, and Slippage)]{%
    Trajectory accuracy of RicoSLAM,
    SLAM~Toolbox~\cite{macenski-jambrecic:joss:2021},
    and GMapping~\cite{grisetti:tro:2007} on four planar sequences of
    the IILABS~3D~\cite{ribeiro:access:2025} dataset
    (Nav~A~Diff, Nav~A~Omni, Loop, and Slippage).
    Both SLAM~Toolbox~\cite{macenski-jambrecic:joss:2021} and
    GMapping~\cite{grisetti:tro:2007} are evaluated with different splitting
    configurations (translation/rotation thresholds).
    RicoSLAM is evaluated in five modes:
    \emph{scan} (each keyframe represents a scan),
    \emph{scan w/reloc.} (scan-based with relocalization within existing graph),
    \emph{local} (each keyframe represents a local map),
    \emph{local w/upd.} (local maps with information matrix update),
    and \emph{local w/reloc.} (local maps with relocalization).
    Evaluation metrics:
    \glslink{ate}{\glsentryfull{ate}},
    \glslink{rte}{\glsentryfull{rte}}, and
    \glslink{rre}{\glsentryfull{rre}}, in which \gls{rte} and \gls{rre} are
    evaluated over a \SI{10}{\metre} segment length.
    Bold entries indicate the best result of online vs offline methods
    within a tolerance of \SI{0.005}{\metre} for \gls{ate},
    \SI{0.05}{\percent} for \gls{rte}, and
    \SI{0.005}{\degree\per\metre} for and \gls{rre}.
    respectively.
    Underlined entries indicate the best result
    within each method (across its variants).
}
\label{tab:slam:results:iilabs3d}
\begin{tabular}{
    l   % method family      (multirow label)
    l|  % method variant
    % 4 sequences from the IILABS3D dataset
    S[table-format=1.3] S[table-format=1.2] S[table-format=1.3] | % nav A diff
    S[table-format=1.3] S[table-format=1.2] S[table-format=1.3] | % nav A omni
    S[table-format=1.3] S[table-format=1.2] S[table-format=1.3] | % loop
    S[table-format=1.3] S[table-format=1.2] S[table-format=1.3]   % slippage
}
\hline
\multicolumn{1}{c}{} &
\multicolumn{1}{c|}{} &
\multicolumn{3}{c|}{\textbf{Nav~A Diff}} &
\multicolumn{3}{c|}{\textbf{Nav~A Omni}} &
\multicolumn{3}{c|}{\textbf{Loop}} &
\multicolumn{3}{c}{\textbf{Slippage}}\\
\multicolumn{1}{c}{} &
\multicolumn{1}{c|}{} &
\multicolumn{1}{c}{\textbf{ATE}} &
\multicolumn{1}{c}{\textbf{RTE}} &
\multicolumn{1}{c|}{\textbf{RRE}} &
\multicolumn{1}{c}{\textbf{ATE}} &
\multicolumn{1}{c}{\textbf{RTE}} &
\multicolumn{1}{c|}{\textbf{RRE}} &
\multicolumn{1}{c}{\textbf{ATE}} &
\multicolumn{1}{c}{\textbf{RTE}} &
\multicolumn{1}{c|}{\textbf{RRE}} &
\multicolumn{1}{c}{\textbf{ATE}} &
\multicolumn{1}{c}{\textbf{RTE}} &
\multicolumn{1}{c}{\textbf{RRE}}\\
\multicolumn{1}{l}{\textbf{Method}} &
\multicolumn{1}{l|}{\textbf{Variant}} &
\multicolumn{1}{c}{(\si{\metre})} &
\multicolumn{1}{c}{(\si{\percent})} &
\multicolumn{1}{c|}{(\si{\degree/\metre})} &
\multicolumn{1}{c}{(\si{\metre})} &
\multicolumn{1}{c}{(\si{\percent})} &
\multicolumn{1}{c|}{(\si{\degree/\metre})} &
\multicolumn{1}{c}{(\si{\metre})} &
\multicolumn{1}{c}{(\si{\percent})} &
\multicolumn{1}{c|}{(\si{\degree/\metre})} &
\multicolumn{1}{c}{(\si{\metre})} &
\multicolumn{1}{c}{(\si{\percent})} &
\multicolumn{1}{c}{(\si{\degree/\metre})}\\
\hline

% ==============================================================================
% GMAPPING
% ==============================================================================

\multirow{3}{*}{GMapping} & 0.3m 0.3rad &
0.138 & \U{1.76} & \U{0.109} &  % nav A diff
0.134 &    2.13  & \U{0.104} &  % nav A omni
0.128 &    3.50  & \U{0.137} &  % loop
0.049 & \U{1.95} & \U{0.118} \\ % slippage
& 0.5m 0.5rad &
   0.135  &    1.88  & 0.121 &  % nav A diff
\U{0.121} & \U{1.89} & 0.128 &  % nav A omni
   0.108  & \U{2.54} & 0.148 &  % loop
\U{0.041} &    2.11  & 0.176 \\ % slippage
& 1.0m 1.0rad &
\U{0.115} & 2.01 & 0.134 &  % nav A diff
   0.123  & 2.94 & 0.243 &  % nav A omni
\U{0.107} & 3.94 & 0.281 &  % loop
   0.060  & 6.53 & 0.540 \\ % slippage

\hline

% ==============================================================================
% SLAM TOOLBOX
% ==============================================================================

\multirow{2}{*}{SLAM Toolbox} & 0.3m 0.3rad &
   0.075  & \UB{0.86} & \U{0.084} &  % nav A diff
   0.081  &  \U{1.40} & \U{0.105} &  % nav A omni
   0.068  &     1.69  &    0.061  &  % loop
\U{0.042} &  \U{5.13} & \U{0.326} \\ % slippage
& 0.5m 0.5rad &
\U{0.053} &    0.94  &     0.121  &  % nav A diff
\U{0.061} &    1.89  &     0.175  &  % nav A omni
\U{0.040} & \U{1.68} & \UB{0.055} &  % loop
   0.064  &    6.00  &     0.487  \\ % slippage

\hline

% ==============================================================================
% RICOSLAM
% ==============================================================================

\multirow{5}{*}{RicoSLAM} & scan &
   0.063 & \UB{0.81} & \B 0.076 &  % nav A diff
   0.053 &     0.86  &    0.081 &  % nav A omni
\B 0.037 &  \B 0.96  &    0.077 &  % loop
\B 0.015 & \UB{0.53} &    0.067 \\ % slippage
& scan w/reloc. &
    0.107  &    1.07 &  \B 0.075  &  % nav A diff
    0.057  &    0.95 &     0.074  &  % nav A omni
    0.053  &    1.36 &     0.081  &  % loop
\UB{0.014} & \B 0.57 & \UB{0.057} \\ % slippage
& local &
\UB{0.047} &    0.91 & 0.084 &  % nav A diff
    0.043  &    0.88 & 0.066 &  % nav A omni
 \B 0.037  & \B 0.96 & 0.081 &  % loop
 \B 0.019  &    0.96 & 0.103 \\ % slippage
& local w/upd. &
    0.058  &     0.92  & 0.084 &  % nav A diff
\UB{0.030} &     0.79  & 0.071 &  % nav A omni
\UB{0.033} & \UB{0.94} & 0.091 &  % loop
 \B 0.018  &     0.95  & 0.111 \\ % slippage
& local w/reloc. &
   0.077 &     0.89  & \UB{0.072} &  % nav A diff
   0.059 & \UB{0.71} & \UB{0.057} &  % nav A omni
   0.042 &     1.11  &  \U{0.068} &  % loop
\B 0.019 &     1.12  &     0.095  \\ % slippage

\hline
\noalign{\vskip 2pt}
\hline

% ==============================================================================
% RICOSLAM (WITH OFFLINE POSE CORRECTION, CONSIDERING LOOP CLOSURES)
% ==============================================================================

\multirow{5}{*}{RicoSLAM$^{\dagger}$} & scan &
\UB{0.022} & \UB{0.74} &     0.066  &  % nav A diff
\UB{0.023} &  \B 0.72  & \UB{0.063} &  % nav A omni
\UB{0.023} &     1.18  &     0.077  &  % loop
 \B 0.016  &     0.64  &     0.076  \\ % slippage
& scan w/reloc. &
    0.041  &     0.89  &     0.070  &  % nav A diff
    0.041  &     0.85  &     0.073  &  % nav A omni
 \B 0.027  &     1.05  & \UB{0.058} &  % loop
\UB{0.014} & \UB{0.57} & \UB{0.057} \\ % slippage
& local &
   0.035 &     0.84  & 0.065 &  % nav A diff
   0.032 &     0.86  & 0.069 &  % nav A omni
   0.030 & \UB{0.99} & 0.067 &  % loop
\B 0.017 &     0.98  & 0.114 \\ % slippage
& local w/upd. &
\B 0.027 &     0.80  & \UB{0.059} &  % nav A diff
\B 0.026 &     0.84  &  \B 0.066  &  % nav A omni
\B 0.027 & \UB{0.99} &     0.070  &  % loop
\B 0.016 &     0.92  &     0.110  \\ % slippage
& local w/reloc. &
   0.043 &     0.85  &    0.072 &  % nav A diff
   0.031 & \UB{0.71} & \B 0.068 &  % nav A omni
   0.039 &     1.06  & \B 0.062 &  % loop
\B 0.019 &     1.12  &    0.095 \\ % slippage

\hline

\multicolumn{14}{l}{%
    \footnotesize
    $^{\dagger}$ offline corrected pose stream
    (see~\secref{sec:slam:online:pose-history}) for the same configurations as
    the online RicoSLAM results, in which}\\
\multicolumn{14}{l}{%
    \footnotesize
    \phantom{$^{\dagger}$} past poses are reconstructed
    \emph{a posteriori} from the optimized keyframe estimates after
    loop closures.}\\


\end{tabular}
\end{table}}


Algorithms:
\begin{itemize}
\item RicoSLAM (online + offline corrected pose, 0.03m resolution,
      voxelization enabled, bi-linear interpolation disabled, M1 representation)
\item SLAM Toolbox
\item GMapping
\item Cartographer (optional, SLAM Toolbox already provides a graph-based formulation)
\item HectorSLAM (optional, very low priority)
\end{itemize}

Metrics:
\begin{itemize}
\item ATE (Absolute Trajectory Error)
\item RTE (Relative Translation Error)
\item RRE (Relative Rotational Error)
\end{itemize}


SLAM Toolbox with 1.00m/1.00rad, 1.50m/1.50rad, and 2.50m (setting minimum angle to 10.0 to avoid splitting due to angular displacement)
fails. Only with 0.3m/0.3rad and 0.5m/0.5rad...

\subsection{Scalability and graph growth with relocalization}\label{sec:slam:results:relocalization}

\todo{TBC}

{\setlength{\tabcolsep}{2pt}
\begin{table}[!t]
\centering
\caption[%
    TBC]{%
    \#n number nodes; \#s number of scans; \#e number of edges%
}
\label{tab:slam:results:iilabs3d:graph}
\begin{tabular}{
    l   % method family      (multirow label)
    l|  % method variant
    % 4 sequences from the IILABS3D dataset
    S[table-format=3.0] S[table-format=3.0] S[table-format=4.0] | % nav A diff
    S[table-format=3.0] S[table-format=3.0] S[table-format=3.0] | % nav A omni
    S[table-format=3.0] S[table-format=3.0] S[table-format=4.0] | % loop
    S[table-format=3.0] S[table-format=3.0] S[table-format=3.0]   % slippage
}
\hline
\multicolumn{1}{l}{} &
\multicolumn{1}{l|}{} &
\multicolumn{3}{c|}{\textbf{Nav~A Diff}} &
\multicolumn{3}{c|}{\textbf{Nav~A Omni}} &
\multicolumn{3}{c|}{\textbf{Loop}} &
\multicolumn{3}{c}{\textbf{Slippage}}\\
\multicolumn{1}{l}{\textbf{Method}} &
\multicolumn{1}{l|}{\textbf{Variant}} &
\multicolumn{1}{c}{\textbf{\#n}} &
\multicolumn{1}{c}{\textbf{\#s}} &
\multicolumn{1}{c|}{\textbf{\#e}} &
\multicolumn{1}{c}{\textbf{\#n}} &
\multicolumn{1}{c}{\textbf{\#s}} &
\multicolumn{1}{c|}{\textbf{\#e}} &
\multicolumn{1}{c}{\textbf{\#n}} &
\multicolumn{1}{c}{\textbf{\#s}} &
\multicolumn{1}{c|}{\textbf{\#e}} &
\multicolumn{1}{c}{\textbf{\#n}} &
\multicolumn{1}{c}{\textbf{\#s}} &
\multicolumn{1}{c}{\textbf{\#e}}\\
\hline

% ==============================================================================
% GMAPPING
% ==============================================================================

\multirow{3}{*}{GMapping} & 0.3m 0.3rad &
{--} & 897 & {--} &  % nav A diff
{--} & 373 & {--} &  % nav A omni
{--} & 757 & {--} &  % loop
{--} & 141 & {--} \\ % slippage
& 0.5m 0.5rad &
{--} & 544 & {--} &  % nav A diff
{--} & 226 & {--} &  % nav A omni
{--} & 459 & {--} &  % loop
{--} &  84 & {--} \\ % slippage
& 1.0m 1.0rad &
{--} & \U{273} & {--} &  % nav A diff
{--} & \U{113} & {--} &  % nav A omni
{--} & \U{231} & {--} &  % loop
{--} &  \U{42} & {--} \\ % slippage

\hline

% ==============================================================================
% SLAM TOOLBOX
% ==============================================================================

\multirow{2}{*}{SLAM Toolbox} & 0.3m 0.3rad &
541 & 541 & 2488 &  % nav A diff
222 & 222 &  409 &  % nav A omni
458 & 458 & 1291 &  % loop
 77 &  77 &  111 \\ % slippage
& 0.5m 0.5rad &
\U{310} & \U{310} & \U{470} &  % nav A diff
\U{126} & \U{126} & \U{163} &  % nav A omni
\U{261} & \U{261} & \U{393} &  % loop
 \U{46} &  \U{46} &  \U{59} \\ % slippage

\hline

% ==============================================================================
% RICOSLAM
% ==============================================================================

\multirow{5}{*}{RicoSLAM} & scan &
     75 &      75 &     168 &  % nav A diff
     31 &      31 &      47 &  % nav A omni
     83 &      83 &     158 &  % loop
     10 &      10 &      11 \\ % slippage
& scan w/reloc. &
     34 & \UB{34} &      36 &  % nav A diff
     18 & \UB{18} & \UB{18} &  % nav A omni
     19 & \UB{19} &      19 &  % loop
 \UB{8} &  \UB{8} &  \UB{7} \\ % slippage
& local &
    106 &     373 &     240 &  % nav A diff
     43 &     155 &      66 &  % nav A omni
     89 &     325 &     168 &  % loop
     15 &      47 &      17 \\ % slippage
& local w/upd. &
    106 &     373 &     238 &  % nav A diff
     43 &     155 &      66 &  % nav A omni
     89 &     325 &     168 &  % loop
     15 &      47 &      17 \\ % slippage
& local w/reloc. &
\UB{17} &      89 & \UB{23} &  % nav A diff
\UB{17} &      80 &      23 &  % nav A omni
\UB{18} &      82 & \UB{18} &  % loop
 \UB{8} &      31 &  \UB{7} \\ % slippage

\hline

\end{tabular}
\end{table}}


Algorithms:
\begin{itemize}
\item RicoSLAM (no relocalization versus with relocalization, 0.03m resolution,
      voxelization enabled, bi-linear interpolation disabled, M1 representation)
\item SLAM Toolbox
\end{itemize}

Metrics:
\begin{itemize}
\item number of nodes along time
\item number of edges along time
\end{itemize}

\subsection{Mapping with dynamics removal}\label{sec:slam:results:dynamics}

{\setlength{\tabcolsep}{1.5pt}
\begin{table}[!t]
\centering
\caption[%
    Effect of the proposed offline dynamics removal mechanism on the occupancy
    probability distribution of the 2D occupancy grid maps generated by RicoSLAM]{%
    Effect of the proposed offline dynamics removal mechanism
    (see~\secref{sec:slam:dynamics:offline}) on the occupancy probability
    distribution of the 2D occupancy grid maps generated by RicoSLAM across
    five datasets of 2D laser scan sequences:
    IILABS~3D~\cite{ribeiro:access:2025}, Avantgarde, Cappero,
    Flowbotic~\cite{flowbotic}, and KUKA.
    Two operating modes considered in this experiment:
    \emph{scan} (each keyframe represents a scan)
    and \emph{local} (each keyframe represents a local map).
    For each sequence and mode, two occupancy grid maps are generated at cell
    resolution $\SI{0.03}{\metre}$:
    the \emph{standard} variant rasterizes the retained scans only through the
    Bresenham~\cite{bresenham:ibm:1965,zingl:uas:2016} line algorithm,
    while the \emph{dyn-rmv} variant applies the offline dynamics removal
    mechanism discussed in \secref{sec:slam:dynamics:offline} before rasterization.
    Evaluation metrics:
    $N_{\mathrm{obs.}}$, the number of valid (observed) cells in each map;
    $\Delta N_{\mathrm{obs.}}$, the difference in $N_{\mathrm{obs.}}$ between
    the \emph{dyn-rmv} and the \emph{standard} variants; and the
    peak/middle mass partition
    ($m_{\mathrm{peak}}, m_{\mathrm{mid}}$) of the empirical occupancy
    distribution over the observed cells of each map, formulated in
    \eqref{eq:slam:results:dynamics:metrics:peak-vs-middle}, for two threshold
    values $\tau \in \{0.05, 0.10\}$. By construction,
    $m_{\mathrm{peak}}(\tau)+m_{\mathrm{mid}}(\tau)=1$ for each $\tau$,
    so a sharper occupancy distribution corresponds to a larger
    $m_{\mathrm{peak}}$ and a smaller $m_{\mathrm{mid}}$.
    Bold entries indicate that all results from the \emph{dyn-rmv} variant
    produce a sharper occupancy distribution than the \emph{standard} variant
    for the same sequence, mode, and threshold.%
}
\label{tab:slam:results:dyn-rmv}
\begin{tabular}{
    l   % Dataset  (multirow label)
    l   % sequence (multirow label)
    l|  % method variant
    ccc | % n_valid
    S[table-format=2.2] S[table-format=1.2] |  % tau  5%: standard
    S[table-format=2.2] S[table-format=1.2] | % tau  5%: dynamics rmv
    % S[table-format=2.2] S[table-format=1.2] |  % tau 10%: standard
    % S[table-format=2.2] S[table-format=1.2] | % tau 10%: dynamics rmv
    S[table-format=2.2] S[table-format=1.2] |  % tau 20%: standard
    S[table-format=2.2] S[table-format=1.2]    % tau 20%: dynamics rmv
}
\hline
\multicolumn{1}{c}{} &
\multicolumn{1}{c}{} &
\multicolumn{1}{c|}{} &
\multicolumn{1}{c}{} &
\multicolumn{1}{c}{} &
\multicolumn{1}{c|}{} &
\multicolumn{4}{c|}{$\tau=0.05$} &
\multicolumn{4}{c}{$\tau=0.10$} \\
% \multicolumn{4}{c}{$\tau=0.20$}\\
\multicolumn{1}{c}{} &
\multicolumn{1}{c}{} &
\multicolumn{1}{c|}{} &
\multicolumn{1}{c}{} &
\multicolumn{1}{c}{} &
\multicolumn{1}{c|}{} &
\multicolumn{2}{c|}{\textbf{standard}} &
\multicolumn{2}{c|}{\textbf{dyn-rmv}} &
% \multicolumn{2}{c|}{\textbf{standard}} &
% \multicolumn{2}{c|}{\textbf{dyn-rmv}} &
\multicolumn{2}{c|}{\textbf{standard}} &
\multicolumn{2}{c}{\textbf{dyn-rmv}}\\
\multicolumn{1}{c}{} &
\multicolumn{1}{c}{} &
\multicolumn{1}{c|}{} &
\multicolumn{1}{c}{\textbf{standard}} &
\multicolumn{1}{c}{\textbf{dyn-rmv}} &
\multicolumn{1}{c|}{} &
\multicolumn{1}{c}{$m_\text{peak}$} &
\multicolumn{1}{c|}{$m_\text{mid}$} &
\multicolumn{1}{c}{$m_\text{peak}$} &
\multicolumn{1}{c|}{$m_\text{mid}$} &
% \multicolumn{1}{c}{$m_\text{peak}$} &
% \multicolumn{1}{c|}{$m_\text{mid}$} &
% \multicolumn{1}{c}{$m_\text{peak}$} &
% \multicolumn{1}{c|}{$m_\text{mid}$} &
\multicolumn{1}{c}{$m_\text{peak}$} &
\multicolumn{1}{c|}{$m_\text{mid}$} &
\multicolumn{1}{c}{$m_\text{peak}$} &
\multicolumn{1}{c}{$m_\text{mid}$}\\
\multicolumn{1}{l}{\textbf{Dataset}} &
\multicolumn{1}{l}{\textbf{Sequence}} &
\multicolumn{1}{l|}{\textbf{Met.}} &
\multicolumn{1}{c}{$N_\text{obs.}$} &
\multicolumn{1}{c}{$N_\text{obs.}$} &
\multicolumn{1}{c|}{$\Delta N_\text{obs.}$} &
\multicolumn{1}{c}{(\si{\percent})} &
\multicolumn{1}{c|}{(\si{\percent})} &
\multicolumn{1}{c}{(\si{\percent})} &
\multicolumn{1}{c|}{(\si{\percent})} &
% \multicolumn{1}{c}{(\si{\percent})} &
% \multicolumn{1}{c|}{(\si{\percent})} &
% \multicolumn{1}{c}{(\si{\percent})} &
% \multicolumn{1}{c|}{(\si{\percent})} &
\multicolumn{1}{c}{(\si{\percent})} &
\multicolumn{1}{c|}{(\si{\percent})} &
\multicolumn{1}{c}{(\si{\percent})} &
\multicolumn{1}{c}{(\si{\percent})}\\
\hline



\multirow{8}{*}{IILABS 3D} & \multirow{2}{*}{Nav A Diff}
 & scan & % [scan-based] Map A: shape=(958, 474) #cells=454092 #n_valid=225571 | Map B: shape=(958, 474) #cells=454092 #n_valid=225091
225571 & 225091 & $-480$ &
95.21 & 4.79 &   % tau  5%: standard
\B 96.19 & \B 3.81 &   % tau  5%: dynamics rmv
96.26 & 3.74 &   % tau 10%: standard
\B 96.74 & \B 3.26 \\  % tau 10%: dynamics rmv
% 97.61 & 2.39 &   % tau 20%: standard
% 97.75 & 2.25 \\  % tau 20%: dynamics rmv
&& local & % [local-map] Map A: shape=(961, 478) #cells=459358 #n_valid=238282 | Map B: shape=(961, 478) #cells=459358 #n_valid=238015
238282 & 238015 & $-267$ &
92.91 & 7.09 &   % tau  5%: standard
\B 94.17 & \B 5.83 &   % tau  5%: dynamics rmv
94.70 & 5.30 &   % tau 10%: standard
\B 95.25 & \B 4.75 \\  % tau 10%: dynamics rmv
% 96.64 & 3.36 &   % tau 20%: standard
% 96.78 & 3.22 \\  % tau 20%: dynamics rmv
 & \multirow{2}{*}{Nav A Omni}
 & scan & % [scan-based] Map A: shape=(951, 301) #cells=286251 #n_valid=211680 | Map B: shape=(951, 301) #cells=286251 #n_valid=211027
211680 & 211027 & $-653$ &
96.30 & 3.70 &   % tau  5%: standard
\B 96.99 & \B 3.01 &   % tau  5%: dynamics rmv
96.84 & 3.16 &   % tau 10%: standard
\B 97.26 & \B 2.74 \\  % tau 10%: dynamics rmv
% 97.83 & 2.17 &   % tau 20%: standard
% 98.01 & 1.99 \\  % tau 20%: dynamics rmv
&& local & % [local-map] Map A: shape=(954, 463) #cells=441702 #n_valid=232092 | Map B: shape=(954, 463) #cells=441702 #n_valid=231639
232092 & 231639 & $-453$ &
93.21 & 6.79 &   % tau  5%: standard
\B 94.43 & \B 5.57 &   % tau  5%: dynamics rmv
94.69 & 5.31 &   % tau 10%: standard
\B 95.27 & \B 4.73 \\  % tau 10%: dynamics rmv
% 96.60 & 3.40 &   % tau 20%: standard
% 96.76 & 3.24 \\  % tau 20%: dynamics rmv
 & \multirow{2}{*}{Loop}
 & scan & % [scan-based] Map A: shape=(1311, 452) #cells=592572 #n_valid=324553 | Map B: shape=(1311, 452) #cells=592572 #n_valid=324249
324553 & 324249 & $-304$ &
96.25 & 3.75 &   % tau  5%: standard
\B 96.90 & \B 3.10 &   % tau  5%: dynamics rmv
96.95 & 3.05 &   % tau 10%: standard
\B 97.28 & \B 2.72 \\  % tau 10%: dynamics rmv
% 97.95 & 2.05 &   % tau 20%: standard
% 98.07 & 1.93 \\  % tau 20%: dynamics rmv
&& local & % [local-map] Map A: shape=(1317, 655) #cells=862635  #n_valid=355885 | Map B: shape=(1317, 655) #cells=862635 #n_valid=355165
355885 & 355165 & $-720$ &
94.72 & 5.28 &   % tau  5%: standard
\B 95.52 & \B 4.48 &   % tau  5%: dynamics rmv
95.87 & 4.13 &   % tau 10%: standard
\B 96.25 & \B 3.75 \\  % tau 10%: dynamics rmv
% 97.20 & 2.80 &   % tau 20%: standard
% 97.32 & 2.68 \\  % tau 20%: dynamics rmv
 & \multirow{2}{*}{Slippage}
 & scan & % [scan-based] Map A: shape=(951, 292) #cells=277692 #n_valid=177664 | Map B: shape=(951, 292) #cells=277692 #n_valid=177476
177664 & 177476 & $-188$ &
98.38 & 1.62 &   % tau  5%: standard
\B 98.64 & \B 1.36 &   % tau  5%: dynamics rmv
98.48 & 1.52 &   % tau 10%: standard
\B 98.68 & \B 1.32 \\  % tau 10%: dynamics rmv
% 98.86 & 1.14 &   % tau 20%: standard
% 98.95 & 1.05 \\  % tau 20%: dynamics rmv
&& local & % [local-map] Map A: shape=(953, 306) #cells=291618  #n_valid=196312 | Map B: shape=(953, 306) #cells=291618 #n_valid=196097
196312 & 196097 & $-215$ &
96.27 & 3.73 &   % tau  5%: standard
\B 96.98 & \B 3.02 &   % tau  5%: dynamics rmv
97.08 & 2.92 &   % tau 10%: standard
\B 97.42 & \B 2.58 \\  % tau 10%: dynamics rmv
% 98.14 & 1.86 &   % tau 20%: standard
% 98.24 & 1.76 \\  % tau 20%: dynamics rmv


\hline



\multirow{7}{*}{avantgarde} & run0 & scan & % Map A: shape=(7173, 8262) #cells=59263326 #n_valid=6944503 | Map B: shape=(7173, 8262) #cells=59263326 #n_valid=6885705
6944503 & 6885705 & $-58798$ &
98.90 & 1.10 &   % tau  5%: standard
\B 99.02 & \B 0.98 &   % tau  5%: dynamics rmv
98.98 & 1.02 &   % tau 10%: standard
\B 99.06 & \B 0.94 \\  % tau 10%: dynamics rmv
% 99.23 & 0.77 &   % tau 20%: standard
% 99.27 & 0.73 \\  % tau 20%: dynamics rmv
& run1 & scan & % Map A: shape=(2514, 3943) #cells=9912702 #n_valid=1961496 | Map B: shape=(2514, 3943) #cells=9912702 #n_valid=1925419
1961496 & 1925419 & $-36077$ &
98.99 & 1.01 &   % tau  5%: standard
\B 99.07 & \B 0.93 &   % tau  5%: dynamics rmv
99.08 & 0.92 &   % tau 10%: standard
\B 99.14 & \B 0.86 \\  % tau 10%: dynamics rmv
% 99.31 & 0.69 &   % tau 20%: standard
% 99.34 & 0.66 \\  % tau 20%: dynamics rmv
& run2 & scan & % Map A: shape=(3549, 2466) #cells=8751834 #n_valid=1868482 | Map B: shape=(3549, 2466) #cells=8751834 #n_valid=1846486
1868482 & 1846486 & $-21996$ &
98.81 & 1.19 &   % tau  5%: standard
\B 98.94 & \B 1.06 &   % tau  5%: dynamics rmv
98.91 & 1.09 &   % tau 10%: standard
\B 99.01 & \B 0.99 \\  % tau 10%: dynamics rmv
% 99.19 & 0.81 &   % tau 20%: standard
% 99.24 & 0.76 \\  % tau 20%: dynamics rmv
& run3 & scan & % Map A: shape=(3545, 2329) #cells=8256305 #n_valid=1753355 | Map B: shape=(3545, 2329) #cells=8256305 #n_valid=1733622
1753355 & 1733622 & $-19733$ &
98.86 & 1.14 &   % tau  5%: standard
\B 98.96 & \B 1.04 &   % tau  5%: dynamics rmv
98.94 & 1.06 &   % tau 10%: standard
\B 99.01 & \B 0.99 \\  % tau 10%: dynamics rmv
% 99.19 & 0.81 &   % tau 20%: standard
% 99.23 & 0.77 \\  % tau 20%: dynamics rmv
& run4 & scan & % Map A: shape=(4092, 2291) #cells=9374772 #n_valid=1880880 | Map B: shape=(4092, 2291) #cells=9374772 #n_valid=1862146
1880880 & 1862146 & $-18734$ &
98.71 & 1.29 &   % tau  5%: standard
\B 98.83 & \B 1.17 &   % tau  5%: dynamics rmv
98.82 & 1.18 &   % tau 10%: standard
\B 98.91 & \B 1.09 \\  % tau 10%: dynamics rmv
% 99.12 & 0.88 &   % tau 20%: standard
% 99.17 & 0.83 \\  % tau 20%: dynamics rmv
& run5 & scan & % Map A: shape=(3999, 2448) #cells=9789552 #n_valid=1965889 | Map B: shape=(3999, 2448) #cells=9789552 #n_valid=1948151
1965889 & 1948151 & $-17738$ &
98.49 & 1.51 &   % tau  5%: standard
\B 98.67 & \B 1.33 &   % tau  5%: dynamics rmv
98.66 & 1.34 &   % tau 10%: standard
\B 98.79 & \B 1.21 \\  % tau 10%: dynamics rmv
% 99.01 & 0.99 &   % tau 20%: standard
% 99.08 & 0.92 \\  % tau 20%: dynamics rmv
& run6 & scan & % Map A: shape=(3544, 2438) #cells=8640272  #n_valid=1645006 | Map B: shape=(3544, 2438) #cells=8640272 #n_valid=1627833
1645006 & 1627833 & $-17173$ &
98.68 & 1.32 &   % tau  5%: standard
\B 98.80 & \B 1.20 &   % tau  5%: dynamics rmv
98.81 & 1.19 &   % tau 10%: standard
\B 98.90 & \B 1.10 \\  % tau 10%: dynamics rmv
% 99.11 & 0.89 &   % tau 20%: standard
% 99.16 & 0.84 \\  % tau 20%: dynamics rmv

\hline



\multirow{2}{*}{cappero} && scan & % [scan-based_rmax-25.0m_res-0.03m] Map A: shape=(2043, 3426) #cells=6999318 #n_valid=842204 | Map B: shape=(2043, 3426) #cells=6999318 #n_valid=840532
842204 & 840532 & $-1672$ &
94.13 & 5.87 &   % tau  5%: standard
\B 94.61 & \B 5.39 &   % tau  5%: dynamics rmv
95.24 & 4.76 &   % tau 10%: standard
\B 95.48 & \B 4.52 \\  % tau 10%: dynamics rmv
% 96.78 & 3.22 &   % tau 20%: standard
% 96.89 & 3.11 \\  % tau 20%: dynamics rmv
&& local & % [local-map-wo-factors-update_rmax-25.0m_res-0.03m_trans-2.50m] Map A: shape=(2055, 3427) #cells=7042485 #n_valid=853734 | Map B: shape=(2055, 3427) #cells=7042485 #n_valid=852711
853734 & 852711 & $-1023$ &
93.67 & 6.33 &   % tau  5%: standard
\B 94.17 & \B 5.83 &   % tau  5%: dynamics rmv
94.90 & 5.10 &   % tau 10%: standard
\B 95.17 & \B 4.83 \\  % tau 10%: dynamics rmv
% 96.58 & 3.42 &   % tau 20%: standard
% 96.69 & 3.31 \\  % tau 20%: dynamics rmv

\hline



\multirow{3}{*}{flowbotic} & 20251024 & local & % Map A: shape=(2758, 3565) #cells=9832270 #n_valid=3815218 | Map B: shape=(2758, 3565) #cells=9832270 #n_valid=3769513
3815218 & 3769513 & $-45705$ &
96.81 & 3.19 &   % tau  5%: standard
\B 97.52 & \B 2.48 &   % tau  5%: dynamics rmv
97.45 & 2.55 &   % tau 10%: standard
\B 97.87 & \B 2.13 \\  % tau 10%: dynamics rmv
% 98.31 & 1.69 &   % tau 20%: standard
% 98.48 & 1.52 \\  % tau 20%: dynamics rmv
& 20251027 & local & % Map A: shape=(2692, 3658) #cells=9847336 #n_valid=3116225 | Map B: shape=(2692, 3658) #cells=9847336 #n_valid=3071177
3116225 & 3071177 & $-45048$ &
97.01 & 2.99 &   % tau  5%: standard
\B 97.68 & \B 2.32 &   % tau  5%: dynamics rmv
97.57 & 2.43 &   % tau 10%: standard
\B 97.97 & \B 2.03 \\  % tau 10%: dynamics rmv
% 98.36 & 1.64 &   % tau 20%: standard
% 98.53 & 1.47 \\  % tau 20%: dynamics rmv
& 20251028 & local & % Map A: shape=(2380, 3471) #cells=8260980 #n_valid=3214031 | Map B: shape=(2380, 3471) #cells=8260980 #n_valid=3166414
3214031 & 3166414 & $-47617$ &
96.98 & 3.02 &   % tau  5%: standard
\B 97.59 & \B 2.41 &   % tau  5%: dynamics rmv
97.64 & 2.36 &   % tau 10%: standard
\B 98.02 & \B 1.98 \\  % tau 10%: dynamics rmv
% 98.44 & 1.56 &   % tau 20%: standard
% 98.62 & 1.38 \\  % tau 20%: dynamics rmv

\hline



kuka && local &  % Map A: shape=(5694, 5914) #cells=33674316 #n_valid=11281724 | Map B: shape=(5694, 5914) #cells=33674316 #n_valid=11169607
11281724 & 11169607 & $-112117$ &
96.28 & 3.72 &   % tau  5%: standard
\B 97.06 & \B 2.94 &   % tau  5%: dynamics rmv
96.95 & 3.05 &   % tau 10%: standard
\B 97.42 & \B 2.58 \\  % tau 10%: dynamics rmv
% 97.96 & 2.04 &   % tau 20%: standard
% 98.15 & 1.85 \\  % tau 20%: dynamics rmv

\hline


\end{tabular}
\end{table}}


Algorithms:
\begin{itemize}
\item RicoSLAM (standard, no offline dynamics removal, 0.03m resolution,
      voxelization enabled, bi-linear interpolation disabled, M1 representation)
\item RicoSLAM (with offline dynamics removal, 0.03m resolution,
      voxelization enabled, bi-linear interpolation disabled, M1 representation)
\end{itemize}

Metrics:
\begin{itemize}
\item probability distribution of the occupancy cell values
\end{itemize}

\subsection{Global bundle adjustment with synthetic data}\label{sec:slam:results:gba}

\todo{TBC}

{\setlength{\tabcolsep}{3pt}
\begin{table}[!t]
\centering
\caption[%
    TBC]{%
    TBC}
\label{tab:slam:results:gba}
\begin{tabular}{
    l     % method family  (multirow label)
    l     % level          (M_0 / M_1 / M_2)
    c |   % variant flags  (V - voxelization)
    % 4 scene groups
    S[table-format=1.3] S[table-format=1.2] |   % no noise
    S[table-format=1.3] S[table-format=1.2]     % with noise (std = 0.03m)
}
\hline
\multicolumn{3}{c|}{} &
\multicolumn{2}{c|}{\textbf{No Noise}} &
\multicolumn{2}{c}{\textbf{$\sigma = \SI{0.03}{\metre}$}}\\
% \multicolumn{3}{c|}{} &
% \multicolumn{2}{c|}{} &
% \multicolumn{2}{c}{()}\\
\multicolumn{3}{c|}{} &
\multicolumn{1}{c}{$\varepsilon_t$} &
\multicolumn{1}{c|}{$\varepsilon_\theta$} &
\multicolumn{1}{c}{$\varepsilon_t$} &
\multicolumn{1}{c}{$\varepsilon_\theta$}\\
\multicolumn{1}{l}{\textbf{Meth.}} &
\multicolumn{1}{c}{$\mathbf{\mathcal{M}_l}$} &
\multicolumn{1}{c|}{\textbf{V?}} &
\multicolumn{1}{c}{(m)} &
\multicolumn{1}{c|}{(\textdegree{})} &
\multicolumn{1}{c}{(m)} &
\multicolumn{1}{c}{(\textdegree{})}\\
\hline

% ==============================================================================
% DISTANCE MAP POINT-TO-POINT GBA ERROR FORMULATION
% ==============================================================================

\multirow{4}{*}{$\mathcal{D}_\mathrm{p2p}$} & \multirow{2}{*}{$\mathcal{M}_0$}
  & \xmark &            % wo/voxelization
0.267 & 2.05 &          % no noise
0.187 & 1.51 \\         % with noise
& & \cmark &            % w/voxelization
0.013 & 0.06 &          % no noise
0.009 & 0.06 \\         % with noise
\cline{2-7}
& \multirow{2}{*}{$\mathcal{M}_1$}
  & \xmark &            % wo/voxelization
0.042 & 0.04 &          % no noise
\B 0.003 & \B 0.02 \\   % with noise
& & \cmark &            % w/voxelization
0.046 & 0.04 &          % no noise
0.004 & \B 0.02 \\      % with noise

\hline

% ==============================================================================
% DISTANCE MAP POINT-TO-PLANE GBA ERROR FORMULATION
% ==============================================================================

\multirow{4}{*}{$\mathcal{D}_\mathrm{p2pl}$} & \multirow{2}{*}{$\mathcal{M}_0$}
  & \xmark &            % wo/voxelization
\B 0.000 & \B 0.00 &    % no noise
0.021 & 0.17 \\         % with noise
& & \cmark &            % w/voxelization
\B 0.000 & \B 0.00 &    % no noise
0.009 & 0.05 \\         % with noise
\cline{2-7}
& \multirow{2}{*}{$\mathcal{M}_2$}
  & \xmark &            % wo/voxelization
0.027 & 0.12 &          % no noise
0.004 & 0.03 \\         % with noise
& & \cmark &            % w/voxelization
0.017 & 0.08 &          % no noise
0.005 & 0.03 \\         % with noise

\hline
\end{tabular}
\end{table}}


\section{Summary and Conclusions}\label{sec:slam:conclusions}

TBC
